\section{Anexo V:}
\subsection{Escenarios planteados para pruebas de usabilidad:}
\noindent\textbf{Escenario 1: Definición de privilegios para auditoría}
\begin{itemize}
    \item \textbf{Rol del usuario:} Administrador de la plataforma.
    \item \textbf{Contexto de la situación:} Un equipo de auditores externos supervisará el rendimiento de los sistemas. Por seguridad, no deben tener permisos de edición; su trabajo se limita a consultar procesos y registrar estados de monitorización.
    \item \textbf{Datos para la tarea:}
        \begin{itemize}
            \item \textbf{Nombre del Rol:} \texttt{Observador Externo}
            \item \textbf{Permisos requeridos:} Aquellos que habiliten las funciones \texttt{monitor:listar} y \texttt{monitor:postmonitor}.
        \end{itemize}
    \item \textbf{Meta:} Crear correctamente el nuevo rol con el acceso mínimo necesario.
\end{itemize}


\noindent\textbf{Escenario 2: Alta de personal temporal y asignación de activos}
\begin{itemize}
    \item \textbf{Rol del usuario:} Administrador de la plataforma.
    \item \textbf{Contexto de la situación:} Un nuevo becario se incorpora hoy al equipo por un periodo limitado. Debes integrarlo en la plataforma para que pueda trabajar, asegurando que tenga acceso a los servicios externos y conste como responsable de su hardware asignado.
    \item \textbf{Datos para la tarea:}
        \begin{itemize}
            \item \textbf{Identidad:} \texttt{Jesús Miguel Mota} (\texttt{j.mmota@test.com}).
            \item \textbf{Condición:} Perfil becario con fecha de salida en 3 meses.
            \item \textbf{Vinculación:} Usuario de GitLab \texttt{j.mmota} y rol \texttt{Observador}.
            \item \textbf{Propiedad:} Responsable de la máquina \texttt{SRV-DOCKER-01}.
        \end{itemize}
    \item \textbf{Meta:} Completar el alta de manera correcta vinculando todos los servicios y activos requeridos.
\end{itemize}

\noindent\textbf{Escenario 3: Despliegue de nuevo servidor de base de datos}
\begin{itemize}
    \item \textbf{Rol del usuario:} Administrador de sistemas.
    \item \textbf{Contexto de la situación:} 
    El laboratorio requiere un nuevo servidor dedicado para una base de datos PostgreSQL. Siguiendo la política de seguridad y red del centro, debes dar de alta la máquina en el panel de control asegurando que tenga conectividad en ambos protocolos (IPv4 e IPv6) y que el certificado de seguridad esté correctamente registrado.
    
    \item \textbf{Datos para la tarea:}
        \begin{itemize}
            \item \textbf{Identificadores:} Hostname \texttt{db-medal-01}, S.O. \texttt{Ubuntu 22.04 LTS}.
            \item \textbf{Stack IPv4:} IP Privada \texttt{10.0.1.50}, Gateway \texttt{10.0.1.1}. (Dejar IP Pública vacía).
            \item \textbf{Stack IPv6:} IP Privada \texttt{fd00::50}, IP Pública \texttt{2001:db8::50}, Gateway \texttt{fd00::1}.
            \item \textbf{Seguridad SSL:} Emisor \texttt{Let's Encrypt}, Caducidad \texttt{12/31/2026}.
            \item \textbf{Recursos:} Configurar con \texttt{16 GB RAM}.
        \end{itemize}
    \item \textbf{Meta:} Registro de manera satisfactoria del servidor. 
\end{itemize}


\noindent\textbf{Escenario 4: Despliegue de API de Procesamiento de Lenguaje Natural}
\begin{itemize}
  \item \textbf{Rol del usuario:} Administrador.
    \item \textbf{Contexto de la situación:} 
    Tu equipo ha finalizado el desarrollo de una nueva API para el análisis de historias clínicas. Siguiendo el flujo de trabajo del laboratorio, debes desplegar este servicio en el clúster de computación, asegurándote de que quede vinculado a la solicitud oficial de investigación para el seguimiento de costes y recursos.
    
    \item \textbf{Datos para la tarea:}
        \begin{itemize}
            \item \textbf{Información Básica:} 
                \begin{itemize}
                    \item \textbf{Nombre del Servicio:} \texttt{NLP-Clinical-Service}
                    \item \textbf{Entorno:} \texttt{EN PRODUCCIÓN}
                    \item \textbf{Severidad:} \texttt{ALTO}
                    \item \textbf{Descripción:} API de procesamiento de texto para la extracción de entidades médicas en el proyecto ELADAIS.
                \end{itemize}
            \item \textbf{Infraestructura y Vínculos:}
                \begin{itemize}
                    \item \textbf{Destino:} Seleccionar el servidor \texttt{SRV-GPU-01}.
                    \item \textbf{Vincular Petición:} Localizar y seleccionar la referencia \texttt{IA-RESEARCH}.
                \end{itemize}
            \item \textbf{Configuración Técnica:}
                \begin{itemize}
                    \item \textbf{Network Binding:} Añadir un mapeo de puertos (ej. 8080:80).
                    \item \textbf{Stack Tecnológico:} \texttt{Python / FastAPI / Docker}.
                    \item \textbf{Acceso:} Activar la opción \texttt{PUBLIC ENDPOINT ACCESS}.
                \end{itemize}
        \end{itemize}
    \item \textbf{Meta:} 
    Lanzar el despliegue (\textit{Launch Deployment}) asegurando que el servicio esté correctamente vinculado a la solicitud \texttt{IA-RESEARCH} y expuesto públicamente.
\end{itemize}

\noindent\textbf{Escenario 5: Creación y vinculación de repositorio en GitLab}
\begin{itemize}
    \item \textbf{Rol del usuario:} Investigador.
    \item \textbf{Contexto de la situación:} 
    Vas a comenzar un nuevo desarrollo para la homogeneización de datos clínicos siguiendo el modelo OMOP CDM. Para asegurar el control de versiones y la colaboración con otros miembros del laboratorio, necesitas inicializar un espacio de trabajo en el servidor de GitLab interno añadiendo al usuario creado previamente.
    
    \item \textbf{Datos para la tarea:}
        \begin{itemize}
            \item \textbf{Nombre del Proyecto:} \texttt{omop-data-homogenization}
            \item \textbf{Visibilidad:} \texttt{Privado} Accesible para el usuario creado previamente.
            \item \textbf{Descripción:} Herramienta de ETL para la transformación de datos hospitalarios al estándar OHDSI.
        \end{itemize}
    \item \textbf{Meta:} Creación del proyecto de gitlab que luego será validado en el gitlab.

\end{itemize}

\noindent\textbf{Escenario 6: Verificación del estado de salud del sistema (Healthcheck)}
\begin{itemize}
    \item \textbf{Rol del usuario:} Cualquier usuario.
    \item \textbf{Contexto de la situación:} 
    Has recibido una alerta de un usuario indicando una posible degradación en el rendimiento de la página del laboratorio. Antes de realizar cualquier intervención técnica o reinicio, necesitas verificar el estado en tiempo real de todos los componentes para identificar el punto exacto de fallo.
    
    \item \textbf{Datos para la tarea:}
        \begin{itemize}
            \item \textbf{Objetivo:} Localizar el panel de \texttt{Healthcheck} global.
            \item \textbf{Acción requerida:} Identificar qué servicio específico reporta un estado distinto a \texttt{UP} o \texttt{HEALTHY}.
        \end{itemize}
    \item \textbf{Meta:} 
    Acceder con éxito a la vista de diagnóstico de salud del sistema y confirmar la disponibilidad de los servicios críticos sin necesidad de consultar logs externos.
\end{itemize}

\noindent\textbf{Escenario 8: Configuración de instancia Code-Server para desarrollo web}
\begin{itemize}
    \item \textbf{Rol del usuario:} Estudiante de una ingeniería.
    \item \textbf{Contexto de la situación:} 
    Necesitas trabajar en el desarrollo de una API en Node.js para tu TFG, pero no dispones de suficiente potencia local en tu portátil. Has decidido solicitar una instancia de \texttt{code-server} (VS Code remoto) en la infraestructura del laboratorio para poder programar desde cualquier lugar con una conexión estable y recursos garantizados.
    
    \item \textbf{Datos para la tarea:}
        \begin{itemize}
            \item \textbf{Identificación:} 
                \begin{itemize}
                    \item \textbf{Nombre del Proyecto:} \texttt{TFG-NodeJS-Backend}
                    \item \textbf{Hostname sugerido:} \texttt{vscode-alejandro}
                \end{itemize}
            \item \textbf{Recursos de Hardware:} 
                \begin{itemize}
                    \item \textbf{CPU:} \texttt{4 Cores}
                    \item \textbf{RAM:} \texttt{8 GB}
                    \item \textbf{Almacenamiento:} \texttt{100 GB SSD}
                \end{itemize}
            \item \textbf{Configuración de Entorno:}
                \begin{itemize}
                    \item \textbf{Sistema Operativo:} \texttt{Ubuntu 22.04}
                    \item \textbf{Imagen Docker:} \texttt{codercom/code-server:latest}
                    \item \textbf{Puerto de red:} Mapear el puerto \texttt{8080} para acceso web.
                \end{itemize}
            \item \textbf{Justificación Técnica:} Necesidad de un entorno de desarrollo persistente con alta disponibilidad para pruebas de integración de API.
        \end{itemize}
    \item \textbf{Meta:} 
    Completar el formulario de petición asegurando que la instancia de \texttt{code-server} esté configurada con los recursos de hardware solicitados y sea accesible mediante el puerto correspondiente.
\end{itemize}

\noindent\textbf{Escenario 9: Reserva de nodo GPU para desarrollo crítico}
\begin{itemize}
    \item \textbf{Rol del usuario:} Investigador del proyecto ELADAIS.
    \item \textbf{Contexto de la situación:} 
    El equipo de investigación del proyecto ELADAIS necesita realizar pruebas intensivas de procesamiento sobre el conjunto de datos del servicio de VIH. Para asegurar que los tiempos de respuesta sean óptimos y no haya interferencias de otros procesos, debes reservar la computación en esa semana.
    
    \item \textbf{Datos para la tarea:}
        \begin{itemize}
            \item \textbf{Recurso a reservar:} \texttt{SRV-GPU-01}.
            \item \textbf{Título del evento:} \texttt{Desarrollo API ELADAIS - Servicio VIH}.
            \item \textbf{Temporalidad:} Una semana completa (del 18 al 25 de mayo de 2026).
            \item \textbf{Descripción:} Pruebas de carga y despliegue de microservicios para la analítica de datos de pacientes con VIH dentro del marco del proyecto ELADAIS.
        \end{itemize}
    \item \textbf{Meta:} 
    Crear con éxito el evento en el calendario de recursos del servidor, asegurando que el nodo \texttt{SRV-GPU-01} aparezca como reservado exclusivamente para dicha actividad durante las fechas indicadas.
\end{itemize}
% TODO: PRUEBAS DE USABILIDAD UX/UI.


\noindent\textbf{Escenario 10: Auditoría técnica y resolución de peticiones pendientes}
\begin{itemize}
    \item \textbf{Rol del usuario:} Supervisor del laboratoio.
    \item \textbf{Contexto de la situación:} 
    Se ha recibido una nueva solicitud de recursos para el proyecto ELADAIS. Como responsable de la infraestructura, debes validar si los recursos solicitados (CPU, RAM, GPU) están justificados y si existe disponibilidad en los nodos indicados antes de comprometer el hardware del laboratorio.
    
    \item \textbf{Datos para la tarea:}
        \begin{itemize}
            \item \textbf{Acción inicial:} Acceder al listado de "Peticiones Pendientes".
            \item \textbf{Identificación:} Localizar la petición con la referencia \texttt{IA-RESEARCH-DEEPLEARNING}.
            \item \textbf{Extracción de datos clave:}
                \begin{itemize}
                    \item Verificar si el nodo solicitado (\texttt{SRV-GPU-01}) tiene carga de trabajo esa semana.
                    \item Revisar la coherencia entre la \textbf{Justificación Técnica} y los \textbf{Recursos de Hardware}.
                    \item Comprobar que la imagen Docker especificada es oficial y segura.
                \end{itemize}
        \end{itemize}
    \item \textbf{Meta:} 
    Tras revisar los detalles, emitir un veredicto formal (Aprobar o Denegar) añadiendo un comentario justificativo para el investigador solicitante.
\end{itemize}
